\adnote{This chapter description from the proposal will function as the "roadmap" for the structure and contents of this dissertation chapter.}

%========================
% Chapter 6 — Conclusion (Proposal)
%========================

% \chapter{Conclusion} % uncomment if your class uses \chapter

This chapter consolidates the dissertation’s contributions, clarifies its limits, and outlines concrete next steps. It closes the loop opened in Chapter 1 by showing how digital audio tools can be re-authored from culturally specific premises. Through \textit{Sonic Speculocultural Technopoiesis}, Black sonic epistemologies were make audible and tangible as computational logic; the case studies yielded patterns and operators of practice, and the discussion articulated their theoretical and practical significance.

\subsubsection{Objectives}
\begin{itemize}[nosep]
	\item Summarize findings and articulate contributions across theory, method, and practice.
	\item Acknowledge limits and define actionable next steps for research and dissemination.
	\item Close the argumentative arc by restating how ST reconfigures defaults through culturally specific computational logic.
\end{itemize}

\subsubsection{Intended Sections/Subsections}
\begin{itemize}[nosep]
	\item Summary of Findings
	\item Contributions
		\begin{itemize}[nosep]
			\item Theoretical
			\item Methodological
			\item Practical / Technical
		\end{itemize}
	\item Limitations
	\item Future Directions
		\begin{itemize}[nosep]
			\item Evaluation Studies
			\item Device and Repertoire Expansion
			\item Open Source/Open Library
			\item Pedagogy and Outreach
		\end{itemize}
	\item Concluding Reflection
\end{itemize}


\begin{comment}
\subsubsection{Summary of Findings}
\begin{itemize}
	\item Defaults in audio tooling (global grids, track-first routing, linear transport) function as cultural arguments rather than neutral conveniences.
	\item Applying ST uncovered leverage points and produced concrete modifications (e.g., Layered Time Lattice, Call--Response Trigger Graph, Groove Anchor Transport, Formation Bus Map; Elastic Time Lattice, Orbit Loop, Memory Return, Phase Drift Window).
	\item Embodied evaluation via practice confirmed cultural authenticity and expressive gain; transmissible documentation made reasoning portable.
	\item The ST cycle (Unsettling $\rightarrow$ Remixing $\rightarrow$ Embodying $\rightarrow$ Transmitting) proved repeatable across two distinct cases.
\end{itemize}

\subsubsection{Contributions}
\subsubsection{Theoretical}
\begin{itemize}
	\item Extends post-phenomenological claims about mediation by specifying cultural configuration at the level of interface and algorithm.
	\item Treats the audio apparatus (engine, controller, code, practice) as material-discursive, showing how operators enact diffractive behavior.
	\item Advances a \textit{specific-to-universal} model: start from situated epistemologies to derive general design principles for defaults.
\end{itemize}

\subsubsection{Methodological}
\begin{itemize}
	\item Formalizes ST as a reproducible research-creation method combining critical technical practice, autoethnography, and small-cohort co-design.
	\item Provides an analysis protocol and evaluation rubric centered on cultural authenticity, expressive affordance, friction reduction, transmissibility, and situated validity.
	\item Establishes a pattern-naming and documentation practice (operator recipes, parameter schemas) that couples code with cultural rationale.
\end{itemize}

\subsubsection{Practical / Technical}
\begin{itemize}
	\item Delivers interaction patterns and temporal operators ready to be implemented as DAW plugins, middleware modules, or bespoke interfaces.
	\item Identifies UI affordances for anchors, returns, layered time, formation grouping, and controlled drift.
\end{itemize}

\subsubsection{Limitations}
\begin{itemize}
	\item \textbf{Scope and cohort.} Small, practice-centered cohort limits generalizability; findings are intentionally situated.
	\item \textbf{Tooling constraints.} Some platforms limit timing and routing control; onboarding is required for elastic or layered time.
	\item \textbf{Ecological validity.} Room acoustics and device diversity were limited; repertoire breadth was modest.
\end{itemize}

\subsubsection{Future Directions}

\subsubsection{Evaluation Studies}
Design mixed-method studies on entrainment, phrasing, perceived expressivity, and error tolerance under different operator settings; compare against baseline tools.

\subsubsection{Device and Repertoire Expansion}
Test across controllers (pads, keyboards, motion sensors), ensembles, and repertoires; stress-test phase drift and layered time under varied tempi and textures.

\subsubsection{Cross-cultural Adaptation}
Collaborate with practitioners from other epistemic traditions to adapt the method while preserving specificity, authorship, and ethical transmission.

\subsubsection{Open Pattern Library}
Maintain a public repository of patterns, operator recipes, short demos, and parameter schemas; support contributions and issue-based refinement.

\subsubsection{Pedagogy and Outreach}
Integrate ST assignments in HCI/interaction design and music technology courses; create tutorial modules on reading and refactoring defaults.

\subsubsection{Concluding Reflection}
The dissertation began by questioning neutrality and ends with a working method for redesigning defaults. The named patterns and operators constitute a portable vocabulary for culturally grounded computation. More broadly, ST demonstrates that re-authoring tools from specific epistemologies can produce general design principles without erasing their origins.


% --- Chapter 6 Key References (keyword-based filter) ---
\begin{refsection}[references.bib] % ensure your ch.6 items have keyword=ch6
	\nocite{*}
	\printbibliography[keyword=ch6,heading=subbibliography,title={Key References}]
\end{refsection}
\end{comment}